\begin{filecontents*}{references.bib}
@article{rmus2025generating,
  title={Generating Computational Cognitive Models using Large Language Models},
  author={Rmus, Milena and Mathony, Marvin and Jagadish, Akshay K. and Ludwig, Tobias and Schulz, Eric},
  journal={arXiv preprint arXiv:2502.00879},
  year={2025}
}
@article{castro2025discovering,
  title={Discovering symbolic cognitive models from human and animal behavior},
  author={Castro, Pablo Samuel and Tomasev, Nenad and Anand, Ankit and Sharma, Navodita and others},
  journal={bioRxiv},
  year={2025},
  publisher={Cold Spring Harbor Laboratory}
}
@article{lueckmann2026discovering,
  title={Discovering mechanistic models of neural activity: System identification in an in silico zebrafish},
  author={Lueckmann, Jan-Matthis and Jain, Viren and Januszewski, Micha{\l}},
  journal={arXiv preprint arXiv:2602.04492},
  year={2026}
}
@article{tilbury2025ai,
  title={AI-discovered tuning laws explain neuronal population code geometry},
  author={Tilbury, Reilly and Kwon, Dabin and Haydaro{\u{g}}lu, Ali and Ratliff, Jacob and others},
  journal={bioRxiv},
  year={2025},
  publisher={Cold Spring Harbor Laboratory}
}
@article{luo2025large,
  title={Large language models surpass human experts in predicting neuroscience results},
  author={Luo, Xiaoliang and Rechardt, Akilles and Sun, Guangzhi and Nejad, Kevin K. and others},
  journal={Nature Human Behaviour},
  year={2025},
  publisher={Nature Publishing Group}
}
@article{luppi2025cognitive,
  title={Cognitive cartography of mammalian brains using meta-analysis of AI experts},
  author={Luppi, Andrea I. and Ali, Hana and Liu, Zhen-Qi and Milisav, Filip and others},
  journal={bioRxiv},
  year={2025},
  publisher={Cold Spring Harbor Laboratory}
}
@article{ghayem2026neurocontext,
  title={NeuroConText: Contrastive learning for neuroscience meta-analysis with rich text representation},
  author={Ghayem, Fateme and Meudec, Rapha{\"e}l and Dock{\`e}s, J{\'e}r{\^o}me and Thirion, Bertrand and Wassermann, Demian},
  journal={Imaging Neuroscience},
  year={2026},
  publisher={MIT Press}
}
@article{wang2025foundation,
  title={Foundation model of neural activity predicts response to new stimulus types},
  author={Wang, Eric Y. and Fahey, Paul G. and Ding, Zhuokun and Papadopoulos, Stelios and others},
  journal={Nature},
  year={2025},
  publisher={Nature Publishing Group}
}
@article{ryoo2025generalizable,
  title={Generalizable, real-time neural decoding with hybrid state-space models},
  author={Ryoo, Avery Hee-Woon and Krishna, Nanda H. and Mao, Ximeng and Azabou, Mehdi and others},
  journal={arXiv preprint arXiv:2506.05320},
  year={2025}
}
@article{wairagkar2025instantaneous,
  title={An instantaneous voice-synthesis neuroprosthesis},
  author={Wairagkar, Maitreyee and Card, Nicholas S. and Singer-Clark, Tyler and Hou, Xianda and others},
  journal={Nature},
  year={2025},
  publisher={Nature Publishing Group}
}
@article{collins2012cognitive,
  title={Cognitive control over learning: Creating, clustering, and generalizing task-set structure},
  author={Collins, Anne GE and Frank, Michael J},
  journal={Psychological review},
  volume={119},
  number={4},
  pages={890},
  year={2012},
  publisher={American Psychological Association}
}
@article{meta2024llama3,
  title={The Llama 3 Herd of Models},
  author={Dubey, Abhimanyu and Jauhri, Abhinav and Pandey, Abhinav and Kadian, Abhishek and others},
  journal={arXiv preprint arXiv:2407.21783},
  year={2024}
}
@article{guo2025deepseek,
  title={DeepSeek-R1: Incentivizing Reasoning Capability in LLMs via Reinforcement Learning},
  author={Guo, Daya and Yang, Dejian and Zhang, Haowei and Song, Junxiao and Zhang, Ruoyu and Xu, Runxin and others},
  journal={arXiv preprint arXiv:2501.12948},
  year={2025}
}
@article{yang2024qwen25,
  title={Qwen2. 5 technical report},
  author={Yang, An and Yang, Baosong and Hui, Binyuan and Zheng, Bo and Yu, Bowen and others},
  journal={arXiv preprint arXiv:2412.15115},
  year={2024}
}

@article{collins2017workingmemory,
  title={Working memory load strengthens reward prediction errors},
  author={Collins, Anne G. E. and Ciullo, Brittany and Frank, Michael J. and Badre, David},
  journal={The Journal of Neuroscience},
  volume={37},
  number={16},
  pages={4332--4342},
  year={2017},
  doi={10.1523/JNEUROSCI.2700-16.2017}
}
@book{kaplan1964conduct,
  title={The Conduct of Inquiry: Methodology for Behavioral Science},
  author={Kaplan, Abraham},
  year={1964},
  publisher={Chandler}
}
\end{filecontents*}

\documentclass[12pt, a4paper]{article}

\usepackage[margin=2.5cm]{geometry}
\usepackage{hyperref}
\usepackage{booktabs}
\usepackage{array}
\usepackage{parskip}
\usepackage{titlesec}
\usepackage{xcolor}
\usepackage[T1]{fontenc}
\usepackage[utf8]{inputenc}
\usepackage[round]{natbib} % Added for proper citation formatting

\hypersetup{
    colorlinks=true,
    linkcolor=blue!60!black,
    urlcolor=blue!60!black,
    citecolor=blue!60!black
}

\titleformat{\section}{\large\bfseries}{}{0em}{}[\titlerule]
\titleformat{\subsection}{\normalsize\bfseries}{}{0em}{}

\title{\textbf{Seminar Thesis on \textit{Generating Computational Cognitive Models
using Large Language Models}}\\
\large Submitted by Christian Block and YingZhu Chen}

\author{}
\date{}

\begin{document}

\maketitle
\tableofcontents

% ─────────────────────────────────────────────────────────────
\section{Executive Summary}

The pipeline proposed in \textit{Generating Computational Cognitive Models
using Large Language Models} (GeCCo) uses Large Language Models (LLMs) to automatically generate, fit, and iteratively refine computational cognitive models as Python functions. Given a task description, participant behavioural data, and a template function, the pipeline in GeCCo prompts an LLM to propose candidate models, fits those proposals to held-out data, and refines them based on feedback derived from predictive performance.

The approach was benchmarked across four cognitive domains---decision making, learning, planning, and memory---using three open-source LLMs of varying parameter size, capacity, and training procedure. Across the four benchmark domains, the strongest LLM-generated models achieved lower (better) BIC scores than the hand-crafted baselines in all four cases, and this advantage was confirmed by a significance test on BIC in decision making ($p<.001$) and working memory ($p<.01$). In learning and planning, the BIC advantage was smaller and not formally tested for significance in the original paper, though the Exceedance Probability strongly favoured the LLM-generated model in learning (0.97 vs.\ 0.03) and more moderately in planning (0.85 vs.\ 0.15). The overall result is therefore best described as broadly competitive to superior, with the strength of evidence varying by domain rather than being uniform.

% ─────────────────────────────────────────────────────────────
\section{Motivation}

Computational cognitive models provide a formal framework for understanding human behaviour. In practice, researchers hand-craft these models, fit them to behavioural data, and iteratively refine them. This process is slow, skill-intensive, and systematically biased due to a person's past exposure to the cognitive literature. Translating a psychological theory into a working algorithm requires lengthy literature reviews and repeated coding cycles. Researchers must simultaneously act as domain experts in cognitive science, theoreticians for data modeling, and software engineers for implementation. Perhaps most importantly, researchers tend to explore only the model classes they already know and favour, potentially missing alternative explanations for the data. This resembles the \emph{streetlight effect}, or what \citet{kaplan1964conduct} called the principle of the drunkard's search: investigation is pulled toward the places where existing tools and theories make searching easiest. The aforementioned limitations narrow the space of models considered, motivating an automated approach that is less tied to any single school of thought in theoretical cognitive science.\\

Three specific LLM capabilities make them well-suited for proposing cognitive models for experimental data.\footnote{The models used were: Llama-3.1-Instruct-70B \citep{meta2024llama3}, DeepSeek-R1-Distill-Llama-3.1-70B \citep{guo2025deepseek}, and Qwen2.5-72B-Instruct \citep{yang2024qwen25}.} First, LLMs can process task descriptions and behavioural data written in natural language, which makes them highly adaptable across different experimental paradigms without the need to provide a custom interface. Second, using broad representational and reasoning abilities, they can identify structure in complex behavioural data and generate plausible psychological hypotheses by reasoning over task structure and behavioural patterns. Third, LLMs can translate those hypotheses directly into executable Python functions. These three abilities have already been utilized in automated statistical modelling through probabilistic program generation, and GeCCo applies the same logic to cognitive science.

% ─────────────────────────────────────────────────────────────
\section{GeCCo's Proposed Pipeline}

Each pipeline run consists of 10 sampling iterations, repeated across 5 independent runs per aforementioned cognitive domain. The process begins by providing the base LLM with a fixed prompt containing a natural language task description, a subset of behavioral data, strict instructions, a domain-specific code template, and starting from the second iteration onwards, iterative feedback. Using this information, the LLM generates three unique cognitive models, ensuring that each features a non-overlapping set of parameters. Each full run took at most eight hours on four Nvidia A100 GPUs (40GB each), which is modest compared to the days typically required by search-based or evolutionary alternatives.\\

Once generated, these models are parsed from the response, converted into executable Python functions, and have their specific parameter names extracted. They are then fitted to a held-out validation dataset that was deliberately excluded from the initial prompt. To prevent the models from getting stuck in local minima, this fitting process is initialized from 20 random starting points.\\

After the fitting stage, the models are evaluated using the Bayesian Information Criterion (BIC) to identify the single best-performing candidate among the three. Finally, the code for this winning model, along with a list of previously used parameter names, is gathered to act as feedback. This feedback is then injected into the prompt for the next iteration, effectively guiding the LLM's future proposals and preventing it from duplicating its past ideas.

\section{Experiments}

Two complementary metrics are used throughout the evaluation. BIC balances goodness of fit against model complexity by penalising the number of free parameters; when two models fit the data equally well, BIC favours the simpler one with fewer parameters. The Exceedance Probability (EXP) is used in the final evaluation stage and quantifies the posterior probability that a given model provides the most prevalent explanation of observed behaviour across the population, translating raw BIC scores into an interpretable confidence measure. Notably, both metrics are computed offline on held-out data using standard optimisation tools rather than by the LLM itself, so it is the statistical fitting procedure---not the LLM's own judgement---that ultimately determines which candidate model is selected.

\subsection{Experiment 1: Decision Making}

People were asked to choose between two fake products (Option A and Option B). They were given ratings from four different "experts," but they were also told that some experts were more trustworthy (had higher "validity") than others. The baseline was the probabilistic Weighted Additive Decision model, which uses four separate parameters to weight the four features. The Llama-generated model achieved a substantially lower BIC (39.40 vs.\ 51.97; $p < .001$). Rather than requiring four parameters, the LLM reduced the problem to a single \emph{discount factor} that smoothly interpolates between three classic decision strategies---\emph{Take the Best}, \emph{Equal Weighting}, and \emph{Weighted Additive}. This single-parameter arbitration parallels the prior-strength parameter used in Bayesian accounts of heuristics as inference under extreme priors, but the LLM arrived at this parsimonious solution without being given any Bayesian framework to work from.

\subsection{Experiment 2: Learning}

Participants played a "slot machine" style game where they repeatedly chose between two options to win money. In a special version of this game, they didn't just see what they won, but they also saw what the other option would have paid out if they had chosen it. The baseline was a four-learning-rate variant of the Rescorla--Wagner model augmented by a perseveration parameter capturing choice stickiness. The LLM-generated model (BIC 77.97 vs.\ 85.24; EXP = 0.97) replaced the asymmetric learning rates with a clean separation between actual and counterfactual value traces, and replaced the perseveration parameter with a psychologically motivated \emph{forgetting parameter} that decays unchosen-option values over time, mirroring the limits of working memory.

\subsection{Experiment 3: Planning}

This is a two-step puzzle game. Participants pick a "magic carpet," which usually takes them to a specific mountain. Once at the mountain, they ask a genie for gold. However, sometimes a "strong wind" blows the carpet to the wrong mountain. This game tests whether people just blindly repeat actions that win them gold, or if they actually "plan ahead" by understanding the rules of the wind and the carpets. The baseline was a Hybrid model combining model-free and model-based value iteration, mixed by a free weighting parameter. The R1-generated model achieved a lower mean BIC (432.47 vs.\ 437.38; EXP = 0.85), although the BIC difference did not reach conventional significance ($p = .27$). By unifying the model-free/model-based dichotomy into a single system using two transition-dependent learning rates, and by omitting the temporal discounting parameter, the model aligned with recent theoretical calls to move beyond strict dual-system accounts of planning.

\subsection{Experiment 4: Working Memory}

People had to learn through trial-and-error which button to press when shown different shapes on a screen. Sometimes they only had to juggle 3 shapes in their head (low cognitive load), and sometimes they had to juggle 6 shapes (high cognitive load). The baseline was the RL-WM model \citep{collins2012cognitive}, in which a slow RL module and a fast but capacity-limited WM module operate independently, with a load-dependent mixing weight. The Llama-generated model outperformed the baseline substantially (BIC 267.25 vs.\ 275.20; EXP = 0.99). Contrary to the standard assumption that reinforcement learning is immune to cognitive load, the LLM's model assigned separate RL learning rates for low- and high-load conditions (\texttt{lr\_low} / \texttt{lr\_high}), suggesting that load modulates the core learning signal---a conclusion consistent with fMRI evidence that working-memory load can strengthen striatal reward-prediction-error signals \citep{collins2017workingmemory}.

\section{Control Experiments}

A mixed-effects regression predicting BIC from \emph{reasoning ability}, \emph{model family}, and \emph{model size} found that Llama outperformed Qwen ($\beta = 5.624$, $p = .03$) despite Qwen having more parameters, suggesting that training data and protocol matter more than raw scale. Reasoning ability showed only a marginal advantage ($\beta = -0.266$, $p = .07$), falling short of conventional significance.

Systematically removing each prompt component and refitting a mixed-effects model on post-ablation BIC scores revealed that all four components contribute meaningfully, though with different magnitudes. Iterative feedback was the strongest single driver of performance ($\beta = -17.040$, $p < .05$), followed by participant data ($\beta = -12.836$, $p < .05$), task description ($\beta = -9.834$, $p < .05$), and the code template ($\beta = -9.359$, $p < .05$). This pattern suggests that the iterative refinement loop---not merely the quality of the initial prompt---is the largest single contributor to GeCCo's performance relative to a single-shot query, though the authors note that the template itself had the smallest effect of the four components, which is relevant to the template-bias concern discussed below.

The LogProber method was applied to Llama prompts across all four domains. Acceleration scores were well below the contamination threshold of 1.0 (range: 0.017--0.648), suggesting that the LLM is reasoning about novel problems rather than recalling memorised solutions from pretraining, though the highest score (0.648, for the working memory task) is notably closer to the threshold than the others and is not discussed further by the authors.

On synthetic datasets with known generating processes, the LLM successfully reverse-engineered the true model in both learning (Rescorla--Wagner) and decision-making (Take the Best; Tallying) domains, supporting the validity of the approach.

Two further robustness checks reinforce this picture. First, replacing BIC with the Akaike Information Criterion (AIC) as the fitness metric produced closely comparable results in the decision-making and planning domains, indicating that GeCCo's performance is not an artefact of the specific fit statistic used to drive its feedback loop. Second, when the hand-crafted code template was itself replaced with a template generated by the LLM, the resulting models fit the data just as well as those built from the human-authored template. This directly addresses the template-bias concern raised in the ablation study above: while the code template does shape the search space to some degree, the pipeline does not appear to depend on a human-designed template to succeed.

CENTAUR is a large foundation model trained to predict human behaviour across cognitive tasks, serving as a proxy for the ceiling of explainable variance. GeCCo matched or significantly exceeded CENTAUR's negative log-likelihood in three of the four domains: it was statistically indistinguishable from CENTAUR in decision making ($p=.78$) and significantly better in learning ($p<.001$) and memory ($p<.001$). In planning, however, GeCCo's negative log-likelihood was numerically \emph{higher} (worse) than CENTAUR's (214.10 vs.\ 206.00), though the difference was not statistically significant ($p=.18$). The overall picture is therefore one of GeCCo matching a large foundation model's predictive power in most domains while remaining fully interpretable, rather than uniformly surpassing it.

% ─────────────────────────────────────────────────────────────
\section{Discussion}
The GeCCo pipeline successfully demonstrates that open-source LLMs can act as automated research assistants, generating and refining psychological models using standard Python code. Across four distinct areas of human cognition, these AI-generated models were broadly competitive with the best hand-crafted models built by human experts: they achieved lower BIC in all four domains, with the improvement confirmed by a significance test in decision making and working memory, while learning and planning showed smaller or non-significant BIC differences alongside favourable Exceedance Probabilities. Crucially, the AI did not merely reproduce known ideas; it frequently proposed simpler or alternative theoretical explanations. For example, the LLM discovered a single discount factor for decision-making heuristics, cleanly separated experienced from counterfactual outcomes in learning tasks, and suggested that cognitive load directly modulates the reinforcement learning rate in working memory tasks. 

Because human cognitive modeling is often inherently limited by the theoretical assumptions researchers start with, they may overlook simpler explanations. GeCCo partially mitigates this "streetlight effect" by acting as a hypothesis-generation tool that expands the set of possible models, though the final interpretation still relies on human scientific judgment. The control experiments provide converging evidence that this success is driven by the AI's ability to learn from its mistakes through iterative feedback, rather than from data contamination, though it is worth noting that several of these analyses (the feature and ablation regressions in particular) rest on a modest number of independent pipeline runs per domain (five), which limits how strongly the relative importance of each factor can be estimated. Furthermore, by evaluating the generated models on strictly held-out data that the LLM never sees during the prompting phase, the pipeline reduces---though does not eliminate---the risk of overfitting; note also that the posterior predictive checks relied on a separate LLM (ChatGPT) to convert the fitting functions into simulation code, an additional, not fully verified step in the pipeline. The AI's resulting code remained highly interpretable to humans while capturing comparable predictable human behaviour to massive, opaque foundation models like CENTAUR in most, though not all, domains.

A major benefit of GeCCo is its high accessibility; the system works out of the box using standard, open-source AI models without the need for expensive or time-consuming retraining. It utilizes a hybrid approach where the AI acts as an intelligent proposal engine to brainstorm the mathematical structures, but the actual testing of how well those models fit the data is done offline using standard, rigorous statistical tools. Thus, the final model selection is determined by BIC on held-out validation data, while the proposal process remains a hybrid search procedure guided by iterative feedback from previous rounds rather than by unverified AI outputs alone. 

Despite these breakthroughs, the authors outline several limitations that chart the course for future research. One limitation is template bias; the AI is currently provided with a basic, human-made code template to guide its syntax, which might unintentionally restrict its creative exploration to specific model classes---although, as noted above, the authors' own follow-up check using an LLM-generated template found comparable performance, which somewhat mitigates this concern. The system's current feedback mechanism relies purely on predictive performance metrics like the BIC. However, a model with a lower BIC is not automatically the best scientific explanation, and human researchers must still verify if the mechanisms are psychologically meaningful and biologically plausible; the authors suggest that future iterations could incorporate a secondary AI to provide qualitative feedback critiquing theoretical logic or biological plausibility. Beyond the limitations the authors raise explicitly, we would add a related concern: because the generated Python code is well-commented and syntactically clean, it can look more authoritative than it is, and researchers evaluating many candidate models across a 10-iteration sampling loop may be tempted to accept a low-BIC model without scrutinising whether its mechanisms are psychologically plausible. 

Finally, the study is limited in scope because it evaluated four specific, controlled laboratory tasks. Future research must extend this pipeline to more complex areas, such as visual perception and natural language processing, which involve richer data formats. It is also vital to test GeCCo's capacity for cross-task generalization, ensuring that an AI-generated cognitive theory can robustly explain behavior across related tasks rather than just fitting a single dataset. Ultimately, these findings suggest that LLMs have the potential to democratize the highly technical field of computational modeling and significantly accelerate the pace of scientific discovery.

\subsection*{Context Within Related Work}
Within the paper's own literature review, GeCCo is positioned alongside a small but growing set of efforts to use LLMs for automated model discovery in cognitive science, most directly \citet{castro2025discovering}, who use evolutionary search over symbolic expressions to discover cognitive models from human and animal behaviour, and prior work that translates verbal strategy descriptions into executable cognitive-model code. Relative to these, the authors emphasise that GeCCo's contribution is a lightweight, purely in-context pipeline: it requires no fine-tuning and no external search algorithm, converging within hours rather than days by using iterative natural-language feedback instead of an evolutionary search loop. This is a meaningful methodological distinction, since it substantially lowers the barrier to entry for individual labs without dedicated model-discovery infrastructure, though it also means GeCCo's search over the space of possible models is comparatively shallow (10 iterations, three candidates per iteration) relative to evolutionary or gradient-based search methods that can explore thousands of candidates.

More broadly, GeCCo belongs to a wider family of ``self-improving'' AI pipelines that use an LLM to propose, test, and refine artefacts based on feedback from an external evaluator. GeCCo fits this general pattern, but with an important safeguard the authors are explicit about: the acceptance criterion for a candidate model is not the LLM's own judgement of quality, but an external, deterministic statistic (BIC or AIC) computed by standard optimisation software on held-out data. This means the refinement loop does not depend on the LLM correctly judging its own outputs, which is a known failure mode in fully self-referential AI pipelines; the trade-off is that GeCCo's notion of "better" is entirely defined by predictive fit, and, as discussed above, fit alone does not guarantee that a model is psychologically or biologically meaningful.

\subsection{Connection to other Papers in the Seminar}
The work by \citet{rmus2025generating} introduces a framework for generating computational cognitive models using LLMs, moving the field from manually handcrafted theories to automated, AI-driven model discovery based on task instructions and behavioral data. This approach is part of a broader paradigm shift in cognitive science and neuroscience, and its relationship to the other provided papers can be understood through several distinct themes: direct methodological parallels, cross-scale model discovery, the expanding role of LLMs in science, and the bridging of top-down and bottom-up neural research.

\citet{rmus2025generating} shares the most direct methodological and conceptual parallel with the research by \citet{castro2025discovering}. Both papers tackle the complex challenge of discovering symbolic cognitive models directly from human and animal behavior. While traditional cognitive science relies on human intuition to formalize theories, both Rmus and Castro demonstrate that modern AI algorithms can autonomously propose, fit, and evaluate mathematical models that explain learning and decision-making processes. They collectively establish a new standard for automated theory building in behavioral science.

This theme of AI-driven automated discovery extends beyond behavior down to the level of neural circuits, linking \citet{rmus2025generating} to the works of \citet{lueckmann2026discovering} and \citet{tilbury2025ai}. While Rmus focuses on high-level cognitive architectures, Lueckmann and Tilbury utilize AI systems and LLM-based tree searches to discover mechanistic models and tuning laws for neural activity. Together, these papers illustrate a unified methodological trend across scales: whether predicting whole-organism behavior or visual cortical neuronal firing rates, AI is increasingly being used to distill complex biological data into parsimonious, interpretable mathematical equations.

Furthermore, \citet{rmus2025generating} exemplifies the growing utility of LLMs as active scientific reasoning engines, a concept heavily supported by several other papers in this collection. \citet{luo2025large} demonstrate that LLMs can surpass human experts in predicting the outcomes of neuroscience experiments, proving that these models possess a deep, actionable understanding of the domain. Similarly, \citet{luppi2025cognitive} and \citet{ghayem2026neurocontext} utilize LLMs and rich text representations to synthesize vast amounts of complex neuroscience literature, performing meta-analyses and mapping cognitive functions to brain regions. Viewed alongside these works, Rmus et al. represents the next logical step: upgrading LLMs from passive synthesizers of existing literature to active generators of novel scientific hypotheses and code.

Finally, the top-down cognitive modeling approach of \citet{rmus2025generating} provides a necessary complement to the bottom-up neural decoding and foundation models explored in the remaining literature. Papers by \citet{wang2025foundation}, \citet{ryoo2025generalizable}, and \citet{wairagkar2025instantaneous} focus intensely on raw neural data, utilizing foundation models and hybrid state-space models to decode brain activity for neuroprostheses and real-time predictions. While these papers excel at translating neural spikes into variables like speech or stimuli, they often lack explicit cognitive frameworks. \citet{rmus2025generating} provides the computational cognitive architectures that could eventually be mapped onto the neural dynamics decoded by these other models. In summary, Rmus et al. serves as a crucial behavioral anchor in a scientific landscape that is rapidly adopting AI to map the brain, decode its signals, and discover its fundamental laws.

% ─────────────────────────────────────────────────────────────
\newpage
\bibliographystyle{apalike}
\bibliography{references}

\end{document}